\documentclass[twocolumn,10pt]{article}
\usepackage[margin=1.8cm]{geometry}
\usepackage{multicol}
\usepackage{amsmath,amssymb,amsfonts}
\usepackage{graphicx}
\usepackage{xcolor}
\usepackage{booktabs}
\usepackage{hyperref}
\hypersetup{colorlinks=true,linkcolor=blue,citecolor=blue,urlcolor=blue}

\definecolor{bctnavy}{RGB}{11,37,69}
\definecolor{bctgreen}{RGB}{27,94,32}

\begin{document}

\title{BCT Letter 104: The Three-Body Problem from BCT Vacuum Topology\\
Topological Regularisation of Gravitational Chaos and Hopf Charge Conservation}

\author{Michel Robert Cabri\'{e}}
\date{March 2026}



\maketitle
\begin{abstract}
The classical gravitational three-body problem is chaotic: no closed-form solution exists, and 
generic initial conditions lead to eventual ejection on timescales exponentially sensitive to 
initial conditions. We show that the BCT Superfluid Lattice Model introduces three fundamental 
modifications that regularise this chaos. First, BCT lattice incompressibility at $a = \ell_P$ 
imposes a minimum separation $|r_i - r_j|_{\min} = R_{\rm BCT}\,\ell_P = 1.843\,\ell_P$, 
removing the singularities that drive chaotic divergence. Second, gravitational interactions 
mediated by OHC vacuum topology produce a trilinear coupling constant 
$g_3 = \alpha_0^2\,(r_{\rm tet}/r_{\rm oct}) = \alpha_0^2 \times 0.5426 = 2.98\times10^{-5}$, 
absent from all classical formulations. Third, the Hopf charge of the three-body vacuum 
configuration is integer-valued; certain configurations are topologically locked and cannot 
escape to infinity on any physical timescale. We derive the BCT stability criterion for 
hierarchical triple systems and predict (Prediction~\#147) that such systems satisfying the 
BCT Hopf criterion will show anomalously long stability lifetimes compared to classical $N$-body 
predictions, testable with existing Gaia~DR3 data. As a corollary, the Sun--Earth--Moon system 
stability over 4.5~Gyr is explained by topological locking at $H_{\rm system} = 3$.
Zero free parameters.
\end{abstract}



\section{Introduction}

The gravitational three-body problem has resisted analytical solution since Newton. Poincar\'{e} 
proved in 1890 that no additional analytic integrals of motion exist beyond the ten classical 
conservation laws~\cite{poincare1890}. Sundman (1912) provided a convergent series 
solution~\cite{sundman1912}, but it converges so slowly as to be physically useless. Numerical 
integration confirms that generic three-body configurations are chaotic, with Lyapunov times 
many orders of magnitude shorter than the age of the universe.

The deepest pathology is the singularity at $r_{ij} = 0$. Classical gravity provides no 
minimum separation; trajectories can approach arbitrarily close encounters, driving divergences 
in velocity and potential energy that make long-term prediction impossible.

The BCT Superfluid Lattice Model~\cite{bct_l1} changes all three fundamental assumptions 
underlying classical chaos: the minimum scale, the nature of gravitational coupling, and the 
topology of the phase space. We develop each modification and derive their consequences for the 
three-body problem.

\section{BCT Lattice Regularisation}

The BCT vacuum is a body-centred tetragonal lattice of Planck spheres at axial ratio 
$c/a = \sqrt{2}$, with nearest-neighbour distance $\ell_P$~\cite{bct_l1}. Adjacent spheres 
are touching; they cannot overlap. This imposes a hard geometric floor:
%
\begin{equation}
a(t) \geq \ell_P = 1.616\times10^{-35}\,\mathrm{m}
\end{equation}
%
For the three-body problem, this translates to a minimum physical separation:
%
\begin{equation}
|r_i - r_j|_{\min} = R_{\rm BCT}\,\ell_P = 1.843\,\ell_P = 2.977\times10^{-35}\,\mathrm{m}
\end{equation}
%
where $R_{\rm BCT} = 1.843$ is the BCT lattice constant. The modified Friedmann 
equation~\cite{bct_l56}:
%
\begin{equation}
H^2 = \frac{8\pi G}{3}\,\rho\!\left(1 - \frac{\rho}{\rho_P}\right)
\end{equation}
%
implies a modified two-body potential at short range:
%
\begin{equation}
V_{\rm BCT}(r) = -\frac{Gm_1 m_2}{r}\left[1 - \left(\frac{r_{\min}}{r}\right)^3\right]
\end{equation}
%
The singularity at $r = 0$ is replaced by a repulsive barrier at $r = r_{\min}$. The phase 
space trajectory is bounded: no orbit can reach $r = 0$.

\section{The Trilinear OHC Coupling Constant}

In classical gravity, three-body interactions decompose into pairwise terms --- there is no 
genuine three-body force. This follows from the linear superposition principle for Newtonian 
gravity and the equivalence principle in GR.

In BCT, gravity is mediated by OHC vacuum topology --- perturbations to the local Hopf charge 
density $H_{\rm local}$ propagating through the superfluid substrate. For a body of mass $m$ 
at position $\mathbf{r}_0$, the OHC perturbation is:
%
\begin{equation}
\delta H(\mathbf{r}) = \frac{Gm}{c^2}\,\alpha_0\,K_0\!\left(\frac{|\mathbf{r}-\mathbf{r}_0|}{\xi_{\rm OHC}}\right)
\end{equation}
%
where $K_0$ is the modified Bessel function of the second kind, $\xi_{\rm OHC}$ is the OHC 
coherence length, and $\alpha_0 = r_{\rm oct}\,r_{\rm tet}/\pi = 0.007408$ is the BCT fine 
structure constant~\cite{bct_l1}.

For three bodies simultaneously perturbing the same OHC region, the perturbations interfere 
through the Bessel resonance modes of the condensate. This produces a genuine trilinear 
coupling:
%
\begin{equation}
\delta H_{123} = \frac{G^2}{c^4}\,\alpha_0^2\,\frac{r_{\rm tet}}{r_{\rm oct}}\,
m_1 m_2 m_3\,I_3(\mathbf{r}_1,\mathbf{r}_2,\mathbf{r}_3)
\end{equation}
%
where the three-body OHC overlap integral is:
%
\begin{equation}
I_3 = \int K_0\!\left(\frac{|\mathbf{r}-\mathbf{r}_1|}{\xi}\right)
K_0\!\left(\frac{|\mathbf{r}-\mathbf{r}_2|}{\xi}\right)
K_0\!\left(\frac{|\mathbf{r}-\mathbf{r}_3|}{\xi}\right)d^3r
\end{equation}
%
The trilinear coupling constant is therefore:
%
\begin{equation}
\boxed{g_3 = \alpha_0^2\,\frac{r_{\rm tet}}{r_{\rm oct}} = (0.007408)^2 \times 0.5426 
= 2.98\times10^{-5}}
\end{equation}
%
This is a pure number. Zero free parameters. The BCT three-body Hamiltonian is:
%
\begin{equation}
\mathcal{H}_{\rm BCT} = \mathcal{H}_{\rm classical} + V_{\rm BCT}^{\rm short} 
+ g_3\,V_{\rm trilinear}(r_{12},r_{13},r_{23})
\end{equation}
%
where $V_{\rm trilinear}$ involves products of the three pairwise OHC perturbation amplitudes. 
This term has no classical analogue and cannot be derived from any modification of the 
two-body potential.

\section{Topological Locking of Hopf-Stable Configurations}

The most significant BCT modification is topological. In BCT, the vacuum configuration of any 
gravitational system carries a Hopf charge $H_{\rm system} \in \mathbb{Z}$~\cite{bct_jhapp}:
%
\begin{equation}
H_{\rm system} = \sum_i H_i + \sum_{i<j} H_{ij} + H_{123} \quad \in \mathbb{Z}
\end{equation}
%
The topological stability theorem~\cite{bct_l45} states that Hopf states are protected by an 
infinite classical barrier (action $S_{\rm Hopf} = 4\pi^2/\alpha_0 = 5329$). $H_{\rm system}$ 
cannot change continuously. For a configuration with $H_{\rm system} = n$, the topological 
lifetime is:
%
\begin{equation}
\tau_{\rm topo} \sim t_{\rm universe}\,\exp\!\left(\frac{S_{\rm Hopf}}{n}\right)
\end{equation}
%
\textbf{BCT Stability Criterion:} A three-body configuration with $H_{\rm system} = n$ is 
topologically locked against ejection if the minimum energy required for ejection exceeds 
the topological transition action $S_{\rm Hopf}/n$.

For hierarchical triple stellar systems ($H_{\rm system} = 1$ or $2$):
\begin{itemize}
\item $H_{\rm system} = 1$: $\tau \sim 10^{2312}\,t_{\rm universe}$ --- effectively infinite
\item $H_{\rm system} = 2$: $\tau \sim 10^{1156}\,t_{\rm universe}$ --- effectively infinite
\end{itemize}

Classical N-body simulations predict ejection on $10^3$--$10^8$ orbital periods. BCT topological 
locking extends this by $\sim 10^{2300}$ for Hopf-stable configurations.

\subsection{The Sun--Earth--Moon System}

The Sun--Earth--Moon system has survived 4.5~Gyr. Laskar (1994) showed the inner solar system 
is chaotic with Lyapunov time $\sim$5~Myr~\cite{laskar1994}. Yet the Moon remains bound. 

BCT explanation: $H_{\rm system} = 3$ (one Hopf unit per body). Topological lifetime:
%
\begin{equation}
\tau_{\rm Moon} \sim 13.8\,{\rm Gyr}\times\exp(5329/3) \sim 10^{770}\,t_{\rm universe}
\end{equation}
%
The Moon is topologically locked. It is not going anywhere. Not because the orbit is classically 
stable --- it is not --- but because the topological barrier against ejection is geometrically 
infinite on any physical timescale.

\section{Prediction \#147 and Test}

\textbf{Prediction \#147 (BCT Three-Body Stability):}
\textit{Hierarchical triple stellar systems satisfying the BCT Hopf stability criterion 
(inner binary separation $< 0.1$~AU, outer companion in 3:1--10:1 period ratio, total 
$H_{\rm system} = 1$ or $2$) will show statistically anomalous stability lifetimes compared 
to classical $N$-body predictions. Specifically: the Gaia~DR3 catalogue cross-matched against 
classical Monte Carlo stability simulations should find that BCT-Hopf-criterion systems have 
survived $\geq 2\times$ longer than simulations predict at the same confidence level.}

This prediction is testable with existing data. The Gaia~DR3 catalogue contains hundreds of 
thousands of resolved stellar multiples with measured orbital parameters~\cite{gaia_dr3}. 
Classical $N$-body stability lifetimes are computable with standard tools (REBOUND, 
MERCURY). The BCT Hopf criterion selects a well-defined subset. No new observations are 
required.

\begin{table}[h]
\begin{tabular}{@{}lcc@{}}
\toprule
Quantity & BCT value & Notes \\
\midrule
$g_3$ (trilinear coupling) & $2.98\times10^{-5}$ & Zero free params \\
$r_{\min}$ (separation floor) & $1.843\,\ell_P$ & Lattice incompressibility \\
$\tau(H=1)$ & $\sim 10^{2312}\,t_U$ & Topological lock \\
$\tau(H=2)$ & $\sim 10^{1156}\,t_U$ & Topological lock \\
$\tau_{\rm Moon}$ & $\sim 10^{770}\,t_U$ & $H_{\rm system}=3$ \\
\bottomrule
\end{tabular}
\caption{BCT three-body key quantities. $t_U = 13.8$~Gyr.}
\end{table}

\section{BCT SAR Application: Subsurface Mapping of Machu Picchu}

The BCT Bessel ratio SAR methodology --- developed for Victorian goldfields subsurface 
mapping (Prediction~\#144~\cite{bct_l102}) --- applies directly to archaeological 
sites. The BCT ratio:
%
\begin{equation}
\mathcal{R}_{\rm BCT}(x,y) = \frac{J_0(\kappa_0\,\sigma_{VV})}{J_0(\kappa_0\,\sigma_{VH})}
\end{equation}
%
where $\kappa_0 = 3.39$ is the first zero of the Josephson-corrected Bessel function and 
$\sigma_{VV}$, $\sigma_{VH}$ are co- and cross-polarisation Sentinel-1 backscatter 
coefficients, is sensitive to subsurface void structures (chambers, tunnels, hydraulic 
channels) at depths of 40--80~m.

Recent Lidar surveys of Machu Picchu reveal extensive surface structures beneath vegetation 
of unexpected scale. BCT SAR complements Lidar by mapping subsurface void topology from 
satellite, non-invasively. Estimated Sentinel-1 data available: $\sim$600 ascending and 
descending pass scenes (2014--2026). The same BCT Bessel ratio processing used to confirm 
the Barrys Reef gold-quartz anomaly (ground-truthed against known mine locations) would 
map the complete subsurface void network of the Sacred Valley.

Prediction: BCT SAR analysis of Machu Picchu will reveal subsurface gold-bearing quartz 
veins consistent with known Inca gold extraction sites, and will identify previously 
unmapped void structures at 20--60~m depth beneath the main plaza complex.

The BCT Foundation will gift this methodology to UNESCO and the Peruvian Ministry of 
Culture as an immediate first act upon registration. Non-commercial, non-invasive, 
open access.

\section{Conclusion}

The BCT three-body problem differs from its classical counterpart in three fundamental 
ways: lattice incompressibility removes singularities; trilinear OHC coupling $g_3$ 
introduces genuine three-body forces absent from all prior formulations; and Hopf charge 
conservation provides topological locking for certain configurations. The chaos is 
regularised --- not eliminated, but constrained. The chaotic phase space is bounded, 
certain configurations are topologically forbidden from ejection, and the trilinear 
coupling correlates all three bodies through the vacuum substrate simultaneously.

The elusive three-body problem has a BCT answer: the geometry of the vacuum provides 
the missing integral of motion. It is not a conserved function on phase space --- it is 
a topological invariant. It cannot be found by Poincar\'{e}'s methods because it is not 
an analytic function. It is discrete, integer-valued, and geometrically guaranteed.

Zero free parameters.

\section*{Acknowledgments}
The author thanks Zeta Vera (CSSC) for computational verification, Nic for the coffee 
during the derivation, and Nyssa and Tegan for sitting on the relevant pages of 
Sundman (1912) at the critical moment.


\begin{thebibliography}{99}
\bibitem{bct_l1} M.R. Cabri\'{e}, BCT Letter 1 --- Emergent Lorentz Invariance. 
    Zenodo 10.5281/zenodo.18958207 (2026).
\bibitem{bct_l56} M.R. Cabri\'{e}, BCT Letter 56 --- BCT Bounce Cosmology. 
    Zenodo (2026).
\bibitem{bct_jhapp} M.R. Cabri\'{e}, BCT Appendix JH --- The Hopfion Interior and OHC. 
    Zenodo 10.5281/zenodo.18975018 (2026).
\bibitem{bct_l45} M.R. Cabri\'{e}, BCT Letter 45 --- Topological Permanence of OHC States. 
    Zenodo (2026).
\bibitem{bct_l102} M.R. Cabri\'{e}, BCT Letter 102 --- BCT SAR Victorian Goldfields. 
    Zenodo (2026).
\bibitem{poincare1890} H. Poincar\'{e}, Sur le probl\`{e}me des trois corps et les 
    \'{e}quations de la dynamique. Acta Math. \textbf{13}, 1--270 (1890).
\bibitem{sundman1912} K.F. Sundman, M\'{e}moire sur le probl\`{e}me des trois corps. 
    Acta Math. \textbf{36}, 105--179 (1912).
\bibitem{laskar1994} J. Laskar, Large-scale chaos in the solar system. 
    Astron. Astrophys. \textbf{287}, L9--L12 (1994).
\bibitem{gaia_dr3} Gaia Collaboration, Gaia DR3 catalogue. 
    A\&A \textbf{674}, A1 (2023).
\end{thebibliography}

\end{document}
